% 12_own_model.tex - compare with the finished state; the same steps for your own model
\section{Finish}

\step{Compare your folder with the finished state}
\cmd{\$WS/model/finished} holds every file as it is at the end of this workshop (without the three changes of the last
section). Compare:
\begin{shell}
$ diff -rqwB $WS/model/finished $WS/work
\end{shell}
\see only lines \cmd{Only in /home/user/Workshop\_2\_Own\_Groundwater\_Model/work: out} (and \cmd{out\_nogw},
\cmd{out\_ascii}, \cmd{out\_refined} and the other output folders, and \cmd{hydrograph.png}, the picture that
\cmd{plot\_hydrograph.py} drew). \cmd{-wB} leaves spacing and empty lines aside, so a
file you typed or copied from this text is not named for those. A line \cmd{Files \dots differ} names a file that is
not as in the finished state: \cmd{diff -wB} without \cmd{-rq} shows the lines.

\section{Do the same for your own model}
\begin{enumerate}
\item Export the HRU polygons with the HRU numbers of your \cmd{.rvh} (GeoJSON), the river lines with their subbasin, a
  DEM and a raster of the aquifer base. Use longitude and latitude for all of them, or one projection for all of them
  (\cmd{:InputCRS}).
\item Choose aquifer profiles, fill the \cmd{AQUIFER\_PROFILE} column, and define the classes and profiles in the
  \cmd{.rvp}.
\item In the coupled HRUs, send the water that your model sends to deep groundwater into \cmd{GROUNDWATER}. Make sure that
  no process takes water out of \cmd{GROUNDWATER} there.
\item Write the \cmd{.rvg}. Start with cells of a few hundred metres to a few kilometres, and refine later.
\item Run, read \cmd{GWModelSummary.txt}, run the audit and plot. Then calibrate (Workshop~6).
\end{enumerate}
